\documentclass[12pt,a4paper]{report}

% ============================================================
% PACKAGES
% ============================================================
\usepackage[utf8]{inputenc}
\usepackage[T1]{fontenc}
\usepackage{mathptmx}                    % Times New Roman equivalent
\usepackage{textcomp}                    % for \textdegree
\usepackage[a4paper,margin=25mm]{geometry}
\usepackage{setspace}\onehalfspacing
\usepackage{graphicx}
\usepackage{booktabs}
\usepackage{tabularx}
\usepackage{longtable}
\usepackage{multirow}
\usepackage{siunitx}
\usepackage{amsmath,amssymb}
\usepackage[hidelinks]{hyperref}
\usepackage{cleveref}
\usepackage{caption}
\usepackage{subcaption}
\usepackage{float}
\usepackage{enumitem}
\usepackage{pdflscape}                   % for landscape pages (Table 2.1)
\usepackage{adjustbox}

% --- Bibliography (natbib + BibTeX) ---
\usepackage[round,authoryear,sort]{natbib}
\bibliographystyle{apalike}              % Harvard-like style, works with natbib on Overleaf

% ============================================================
% CUSTOM COMMANDS
% ============================================================
\newcommand{\Heworst}{$H_{\mathrm{e,worst}}$}
\newcommand{\Wemaxworst}{$W_{\mathrm{e,max,worst}}$}
\newcommand{\Toppeak}{$T_{\mathrm{op,peak}}$}
\newcommand{\hoursgt}{hours\_gt26\_night}
\newcommand{\heatkwh}{heat\_kWh\_m\textsuperscript{2}}

% ============================================================
\begin{document}

% --- Title Page ---
\begin{titlepage}
\centering
\vspace*{2cm}
{\Large UCL Bartlett School of Environment, Energy and Resources\\[0.5cm]}
{\large BSc SBEER --- BENV0165 Dissertation\\[2cm]}
{\LARGE\bfseries Machine Learning Surrogate Models for Predicting Summertime Overheating Risk in UK Residential Buildings\\[2cm]}
{\large Houlin Bao\\[0.5cm]}
Supervisor: Dr Michael Gutland\\[1cm]
Word count: $\sim$8{,}000\\[2cm]
\today
\end{titlepage}

% --- Abstract ---
\begin{abstract}
\noindent
% TODO: Write abstract (excluded from word count)
\end{abstract}

\tableofcontents
\listoffigures
\listoftables

% ============================================================
% CHAPTER 1: INTRODUCTION
% ============================================================
\chapter{Introduction}
\label{ch:introduction}

\section{Background and Motivation}
\label{sec:background}

Summertime overheating in UK residential buildings has emerged as a significant and growing public health concern. The 2003 European heatwave contributed to over 2{,}000 excess deaths across England and Wales \citep{Johnson2005} --- an event once regarded as a climatic extreme, but whose equivalent severity is projected to become commonplace by the 2040s under medium-high emissions scenarios \citep{Murphy2009}. Since indoor environments mediate a substantial proportion of heat exposure, particularly for vulnerable groups such as the elderly, the thermal performance of the housing stock is directly linked to population health outcomes. \citet{Heaviside2016} found that the Urban Heat Island (UHI) effect alone accounted for approximately half of heat-related mortality during the 2003 event in the West Midlands, underscoring the compound risks facing urban dwellers.

The UK housing stock is particularly ill-suited to these emerging conditions. Approximately 20\% of English dwellings predate 1919, reflecting a construction philosophy optimised for resisting cold and damp winters: solid masonry walls, limited ventilation provision, and no systematic summer shading \citep{Piddington2020}. Successive energy efficiency policies --- from Part~L of the Building Regulations to the Code for Sustainable Homes --- have increased insulation levels and airtightness across the stock, reducing winter heat loss but simultaneously diminishing the capacity for passive summer cooling. \citet{Mulville2016} demonstrated that as energy efficiency standards increased, so did projected overheating risk, particularly under mid-century climate scenarios. This tension between winter energy conservation and summer thermal comfort motivates the development of rapid screening tools capable of evaluating overheating risk across the diverse UK housing stock.


\section{Existing Literature and Research Gap}
\label{sec:lit-gap}

The academic response to residential overheating has grown considerably since the mid-2000s. \citet{Symonds2016} developed a neural-network-based England-wide model linking building characteristics, climate, and urban morphology to indoor thermal exposures; \citet{Botti2022} demonstrated the viability of metamodelling for overheating risk prediction in London apartments. These studies confirm that building characteristics --- including wall type, built form, and floor level --- are strong determinants of overheating risk, in some cases more influential than climatic variability. However, current compliance tools such as EnergyPlus and IES-VE, which underpin assessment frameworks like CIBSE TM52 and TM59, are computationally intensive and require specialist expertise, rendering them poorly suited for rapid, portfolio-level assessment of the existing stock.

Critically, no published surrogate model simultaneously targets UK detached and semi-detached dwellings --- the predominant built forms in the English housing stock --- while parameterising geometry, orientation, and window-to-wall ratio (WWR) as continuous variables and predicting continuous TM52/TM59 overheating metrics alongside heating energy demand.


\section{Research Aim, Objectives, and Methods}
\label{sec:aim}

This dissertation develops a computationally efficient, Python-based surrogate model using supervised machine learning to predict summertime overheating risk in UK detached dwellings, assessed against CIBSE TM52 adaptive comfort thresholds. The specific objectives are:

\begin{enumerate}[label=(\roman*)]
    \item construct a parametric EnergyPlus building energy model for a standard UK detached archetype under current climate conditions;
    \item generate a synthetic training dataset of 2{,}000 simulations through systematic variation of 14 building design parameters using Latin Hypercube Sampling;
    \item train and evaluate three supervised learning architectures --- Random Forest, XGBoost, and Multi-Layer Perceptron --- for predicting five continuous overheating and energy performance targets; and
    \item conduct SHAP-based sensitivity analysis to identify the most influential design parameters and discuss implications for stock-level overheating assessment.
\end{enumerate}


\section{Contribution to Knowledge and Dissertation Structure}
\label{sec:contribution}

This study makes an original contribution by providing the first openly documented surrogate model for overheating risk in UK detached housing that treats geometry, orientation, and WWR as continuous variables --- addressing a gap identified in the existing metamodelling literature. By bridging the gap between high-fidelity dynamic thermal simulation and millisecond-scale machine learning inference, the model offers a practical tool for early-stage design screening and retrofit prioritisation at portfolio scale.

The dissertation proceeds as follows: \Cref{ch:litreview} reviews relevant literature on overheating risk, surrogate modelling, and machine learning in building performance simulation; \Cref{ch:methods} details the methodology; \Cref{ch:results} presents results; \Cref{ch:discussion} discusses findings and limitations; \Cref{ch:conclusion} draws conclusions and recommends future research directions.


% ============================================================
% CHAPTER 2: LITERATURE REVIEW
% ============================================================
\chapter{Literature Review}
\label{ch:litreview}

\section{Surrogate Modelling in Building Performance Simulation}
\label{sec:surrogate-modelling}

Evaluating the thermal performance of a building design requires solving coupled heat and mass balance equations across an annual period using dynamic simulation engines such as EnergyPlus, IES-VE, or TRNSYS \citep{Shatat2015}. A single simulation may take several minutes; parametric studies of thousands of design variants are therefore computationally prohibitive without an alternative approach. Surrogate modelling provides that alternative: a machine learning model is trained on a pre-computed dataset of simulation input--output pairs, enabling millisecond-scale predictions for new parameter combinations without re-running the physics-based model \citep{Westermann2019}. This constitutes a supervised learning problem --- the algorithm learns a mapping from labelled examples (here, EnergyPlus outputs) to generalise across the design space.

The methodology is well-established. \citet{Westermann2019} reviewed over 50 surrogate studies in building design and identified artificial neural networks (ANNs) and Gaussian Processes as the dominant architectures, Latin Hypercube Sampling (LHS) as the standard sampling method, and annual energy demand as the most common target --- with overheating risk comparatively neglected. \citet{Hou2024} subsequently formalised the workflow into a nine-step protocol and used it to audit 68 published studies, finding that the majority lacked adequate reporting of sampling strategy, data preprocessing, and hyperparameter selection. The present study follows this protocol's recommended workflow --- explicitly reporting the sampling method, data preprocessing steps, model architectures, and evaluation metrics --- to ensure methodological transparency and reproducibility.

Several international studies demonstrate the versatility of the approach. \citet{Lawrence2021} trained ANN and multivariate regression surrogates on 6{,}000 EnergyPlus simulations of four Canadian residential archetypes with 19 continuous and 4 categorical parameters, showing that a single model could serve multiple building geometries without loss of accuracy (combined $R^2 = 0.929$). \citet{BagheriEsfeh2025} extended this comparative approach to Random Forest, XGBoost, and feedforward neural networks with SHAP-based feature importance --- a workflow the present study adopts. However, these studies target energy consumption in heating-dominated climates; overheating risk prediction and the UK regulatory context remain largely unaddressed.


\section{Overheating Risk in UK Housing}
\label{sec:overheating}

\subsection{The UK Context and Regulatory Framework}
\label{sec:uk-context}

The UK housing stock is uniquely vulnerable to summertime overheating. Approximately 20\% of English dwellings predate 1919, built with solid masonry walls and limited ventilation for cold winters \citep{Piddington2020}. Successive energy efficiency policies have compounded this vulnerability: the BRE Trust \citeyearpar{BRETrust2020} documented that average SAP ratings rose from 45 to 62 between 1996 and 2017, reflecting improved insulation that simultaneously reduces passive summer cooling capacity. \citet{Shatat2015} demonstrated this trade-off directly, showing that comprehensive retrofit of a Nottingham semi-detached house cut heating energy by 80\% but left summer performance unassessed --- an increasingly important omission given the scale of the semi-detached stock in England.

The regulatory response is codified in CIBSE TM52 \citep{CIBSE2013}, which defines three adaptive comfort criteria: Criterion~1 (Hours of Exceedance $H_e < 3\%$), Criterion~2 (Daily Weighted Exceedance $W_e \leq 6$), and Criterion~3 (absolute upper limit of $T_{\max} + 4$\,K). TM59 \citep{CIBSE2017} extends this to residential buildings, adding a nighttime bedroom criterion (operative temperature $> 26$\,\textdegree C for no more than 32 hours). Compliance assessment requires dynamic thermal simulation under free-running conditions --- a computationally intensive process that, as \citet{Botti2022} observed, is often performed too late in design for passive mitigation strategies to remain feasible. This computational bottleneck motivates the development of surrogate models that can deliver equivalent screening at a fraction of the time.

\subsection{Existing UK Surrogate Models and Their Limitations}
\label{sec:uk-surrogates}

\citet{Symonds2016} developed the first England-wide overheating metamodel: 384 separate neural networks, each trained on 500 LHS samples with 8 continuous inputs. Single-hidden-layer networks with 20 neurons achieved $R^2$ of 0.81--0.97. However, geometry was fixed per archetype and neither orientation nor WWR was parameterised, limiting the model's ability to capture the diversity of the existing stock. \citet{Botti2022} addressed the orientation gap by testing eight discrete orientations for London apartments, with XGBoost achieving $R^2 \approx 0.99$ and identifying window opening factor as the dominant predictor. Yet their model covered only flats with a single Part~L1A-compliant envelope --- a narrow slice of the housing stock. \citet{Petrou2021} investigated wall U-value uncertainty in semi-detached houses but relied entirely on the Symonds metamodel rather than generating new training data.

Across these studies, a consistent pattern emerges: each advances one aspect of the problem while leaving others constrained. Geometry remains fixed, or the building typology is limited to flats, or the parameter space is too narrow to represent the full diversity of the stock. No existing UK surrogate simultaneously parameterises geometry, orientation, and WWR as continuous variables for detached and semi-detached houses --- which together constitute the majority of the English housing stock \citep{EHS2022}.


\section{Machine Learning Methods}
\label{sec:ml-methods}

The three ML architectures employed in the present study each offer distinct advantages for building performance prediction. XGBoost \citep{Chen2016} builds decision trees sequentially, with each new tree trained to correct the residual errors of its predecessors using the gradient of the loss function. This sequential error-correction is particularly effective for building datasets where certain parameter regions --- such as the threshold-like onset of overheating at high WWR --- exhibit highly non-linear responses, as boosted trees can allocate additional model capacity to these difficult regions without overfitting the smooth areas \citep{Botti2022,Shirzadi2025}.

Where XGBoost excels at capturing localised non-linearities, Multi-Layer Perceptron (MLP) neural networks offer smooth function approximation across the entire parameter space. An MLP consists of an input layer, one or more hidden layers with non-linear activation functions (typically ReLU), and an output layer, trained via backpropagation to minimise prediction error \citep{Lawrence2021}. The universal approximation theorem guarantees that a sufficiently wide network can approximate any continuous function, making MLPs inherently suited to the complex, multi-parameter relationships governing building thermal performance. However, this flexibility comes at a cost: MLPs are more susceptible to overfitting and require careful tuning of architecture, learning rate, and regularisation --- typically addressed through early stopping and dropout \citep{Hou2024}.

Random Forest \citep{Breiman2001} serves as a robust baseline, averaging predictions across hundreds of independently trained decision trees to reduce variance. Its built-in Mean Decrease Impurity metric provides computationally cheap feature importance rankings, offering a useful cross-check against the more rigorous but computationally expensive SHAP analysis \citep{Lundberg2017}. SHAP decomposes individual predictions into additive feature contributions grounded in cooperative game theory, enabling both global importance ranking and directional interpretation of each parameter's effect --- critical for translating surrogate model outputs into physically meaningful design guidance \citep{BagheriEsfeh2025}.


\section{Research Gap}
\label{sec:research-gap}

\Cref{tab:lit-comparison} positions the present study against the existing literature. No published surrogate simultaneously: (a)~targets UK detached and semi-detached houses; (b)~parameterises geometry, orientation, and WWR as continuous variables; (c)~predicts continuous TM52/TM59 overheating metrics alongside heating energy; and (d)~provides an openly documented Python implementation. This dissertation addresses this gap with a 14-parameter surrogate trained on 2{,}000 EnergyPlus simulations per archetype, predicting five continuous targets using RF, XGBoost, and MLP with SHAP-based sensitivity analysis. The two archetypes differ in a physically significant way --- the detached house has four exposed facades while the semi-detached shares one party wall, reducing both heat loss and solar gain on that face --- enabling comparison of how built form mediates overheating risk across the same parameter space. The following chapter describes the methodology for generating the training dataset and developing these surrogate models.

% --- Table 2.1: Literature comparison (landscape page) ---
\begin{landscape}
\begin{table}[p]
\centering
\caption{Comparison of surrogate modelling studies for building overheating and energy prediction. Bold entries indicate features unique to the present study.}
\label{tab:lit-comparison}
\footnotesize
\begin{tabularx}{\linewidth}{lllcclllcl}
\toprule
\textbf{Study} & \textbf{Location} & \textbf{Building Typology} & \textbf{Params} & \textbf{Samples} & \textbf{WWR} & \textbf{Orient.} & \textbf{Targets} & \textbf{Best $R^2$} & \textbf{Key Limitation} \\
\midrule
Symonds et al.\ (2016)       & England        & 8 archetypes              & 8+4   & 500/combo   & Fixed   & No        & Overheating + energy        & 0.81--0.97  & Fixed geometry       \\
Botti et al.\ (2022)         & London         & Flats                     & $\sim$5 & 2{,}304   & Fixed   & 8 disc.   & TM59 overheating            & $\approx$0.99 & Flats only        \\
Lawrence et al.\ (2021)      & Toronto        & 4 residential             & 19+4  & 6{,}000     & Sampled & 4 disc.   & Annual energy               & 0.929       & No overheating       \\
Petrou et al.\ (2021)        & England        & Semi-detached             & 13    & 300         & Fixed   & Fixed     & Overheating                 & N/A         & Uses existing model  \\
Bagheri-Esfeh et al.\ (2025) & Iran           & Office                    & 4     & ---         & ---     & Sampled   & Energy, CO\textsubscript{2} & 0.999       & Non-UK; office       \\
Shirzadi et al.\ (2025)      & Toronto        & Midrise apt.              & 8     & 2{,}000     & ---     & Sampled   & EUI + cost                  & $>$0.9      & Non-UK; commercial   \\
Jiang \& Kazmi (2024)        & Netherlands    & Net-zero residential      & Data  & Obs.        & N/A     & N/A       & Overheating (future)        & N/A         & Monitoring data req. \\
\midrule
\textbf{This study}          & \textbf{London} & \textbf{Detached/Semi-D}  & \textbf{14} & \textbf{2{,}000/arch.} & \textbf{10--40\%} & \textbf{0--360\textdegree} & \textbf{TM52/59 + heating} & \textbf{0.955} & \textbf{TMYx climate} \\
\bottomrule
\end{tabularx}
\end{table}
\end{landscape}


% ============================================================
% CHAPTER 3: METHODS
% ============================================================
\chapter{Methods}
\label{ch:methods}

\section{Methodological Overview}
\label{sec:methods-overview}

The methodology follows a three-stage workflow consistent with the surrogate modelling protocol proposed by \citet{Hou2024}: (i)~reference building energy model development and validation in EnergyPlus; (ii)~systematic parameter variation using Latin Hypercube Sampling and parallel batch simulation to generate a synthetic training dataset; and (iii)~surrogate model training, evaluation, and sensitivity analysis. \Cref{fig:workflow} illustrates the end-to-end workflow; \Cref{tab:methodology} summarises the key methodological decisions.

\begin{figure}[htbp]
    \centering
    \includegraphics[width=\textwidth]{figures/workflow.png}
    \caption{Methodology workflow diagram: OpenStudio geometry $\rightarrow$ EnergyPlus IDF template $\rightarrow$ LHS parametric study $\rightarrow$ ML surrogate training $\rightarrow$ SHAP analysis. Adapted from \citet{Lawrence2021}.}
    \label{fig:workflow}
\end{figure}

\begin{table}[htbp]
\centering
\caption{Surrogate model development methodology.}
\label{tab:methodology}
\footnotesize
\begin{tabularx}{\textwidth}{llX}
\toprule
\textbf{Stage} & \textbf{Decision} & \textbf{Choice} \\
\midrule
\multirow{8}{*}{\parbox{3cm}{Dataset\\Development}}
    & Model Intent       & Early-stage overheating screening; sensitivity analysis \\
    & Target Variables   & 5 continuous: $H_{e,\mathrm{worst}}$, $W_{e,\mathrm{max,worst}}$, $T_{\mathrm{op,peak}}$, hours\_gt26\_night, heat\_kWh\_m\textsuperscript{2} \\
    & Archetypes         & 2-storey detached + semi-detached (UK cavity wall) \\
    & Climate            & London TMYx 2009--2023 (WMO 037700) \\
    & Simulation         & EnergyPlus v25.1 \\
    & Base Model         & OpenStudio geometry + BENV0166 physics model \\
    & Parameters         & 14 continuous --- see \Cref{tab:parameters} \\
    & Sampling           & LHS, $n = 2{,}000$ per archetype (seed = 42) \\
\midrule
\multirow{2}{*}{\parbox{3cm}{Data\\Preprocessing}}
    & Train/Val/Test     & 70/15/15, random\_state = 42 \\
    & Feature Engineering & No scaling (tree models); StandardScaler (MLP) \\
\midrule
\multirow{3}{*}{\parbox{3cm}{Surrogate\\Development}}
    & Random Forest      & sklearn, 200 trees, max\_features = $\sqrt{d}$, per-target \\
    & XGBoost            & 400 estimators, lr = 0.05, max\_depth = 6, per-target \\
    & MLP                & PyTorch, 14$\to$128$\to$64$\to$32$\to$5, Adam, early stopping \\
\bottomrule
\end{tabularx}
\end{table}


\section{Simulation Engine and Climate Data}
\label{sec:simulation}

EnergyPlus version 25.1 (US Department of Energy) was selected as the dynamic thermal simulation engine. EnergyPlus is open-source, command-line scriptable, and widely adopted for TM52/TM59 overheating compliance modelling in the UK context \citep{Symonds2016,Botti2022}. Its command-line interface enables automated batch execution from Python, which is essential for running large parametric studies. EnergyPlus includes the Foundation:Kiva two-dimensional finite element ground heat transfer solver, which provides physically accurate ground-coupled heat transfer for parameterised floor geometries --- a capability not available through simplified ground boundary conditions, and adopted here following \citet{Lawrence2021}.

The weather file used is the TMYx 2009--2023 Typical Meteorological Year for London St James's Park (WMO 037700; 51.50\textdegree N, 0.13\textdegree W; elevation 5\,m), obtained from \texttt{climate.onebuilding.org}. TMYx files are constructed from the most recent 15-year observational record using the Sandia method, representing current-climate baseline conditions. The peak summer dry-bulb temperature in this file is 30.8\,\textdegree C. All simulations use a full annual run period (1~January to 31~December) with a 15-minute simulation timestep and hourly output reporting. The use of a TMYx file rather than a CIBSE Design Summer Year (DSY) is a methodological choice discussed in \Cref{sec:limitations}; the surrogate model's architecture permits rapid re-evaluation under alternative weather files without re-running EnergyPlus.


\section{Reference Building Energy Models}
\label{sec:reference-models}

\subsection{Detached Archetype}
\label{sec:detached}

The primary reference building is a two-storey detached house comprising three thermal zones: Space~101 (ground floor living areas), Space~102 (first floor bedrooms), and Space~1 (unoccupied attic beneath a pitched roof). The geometry was created in OpenStudio and exported as an EnergyPlus IDF, with Foundation:Kiva ground heat transfer, WindowMaterial:Shade, Daylighting:Controls, ThermalResilienceSummary outputs, and NECB residential schedules.

Building geometry is parameterised through width (6--10\,m), length (6--12\,m), and storey height (2.4--3.0\,m), following \citet{Lawrence2021}. Orientation is sampled continuously (0--360\textdegree), capturing the facade solar exposure effect identified by \citet{Botti2022}. The external wall follows a UK cavity wall construction: 25\,mm stucco, 200\,mm brick masonry, parameterised insulation ($\lambda = 0.05$\,W/mK, 10--300\,mm), and 19\,mm gypsum board. Roof, slab, and glazing constructions are similarly parameterised, with insulation ranges spanning pre-1919 stock to Passivhaus-equivalent levels, consistent with \citet{Symonds2016} and Part~L \citep{HMGov2021}. Full construction details are provided in Appendix~A.

Glazing uses the SimpleGlazingSystem object with parameterised U-factor (0.80--5.00\,W/m\textsuperscript{2}K) and SHGC (0.20--0.70). Window-to-wall ratio (10--40\%) determines the total glazed area per facade: three windows per floor on the north and south faces, one on each of the east and west faces. Each window includes a Shading:Overhang with sampled projection depth (0--1.0\,m). Heating is provided by IdealLoadsAirSystem at 21\,\textdegree C daytime / 18\,\textdegree C night setback; cooling is fixed at 100\,\textdegree C to ensure free-running summer conditions per TM59 \citep{CIBSE2017}. Natural ventilation uses eight WindandStackOpenArea objects (four orientations $\times$ two zones), triggered when indoor temperature exceeds 24\,\textdegree C and outdoor exceeds 10\,\textdegree C. Infiltration follows the FlowCoefficient method with sampled flow coefficient (0.001--0.012\,m\textsuperscript{3}/s), spanning Passivhaus airtightness to leaky pre-1919 dwellings \citep{Stephen2000}. Occupancy is derived from the SAP 2012 Table~1b formula based on total floor area (2.4--3.8 occupants), maintaining physical consistency between building size and household size.

\begin{figure}[htbp]
    \centering
    \includegraphics[width=0.8\textwidth]{figures/3d_model.png}
    \caption{3D axonometric view of the three-zone detached archetype model.}
    \label{fig:3d-model}
\end{figure}


\subsection{Semi-Detached Archetype}
\label{sec:semi-detached}

The semi-detached model shares the same three-zone structure, construction assemblies, schedules, and parameter ranges as the detached archetype, with three modifications arising from the shared party wall. First, one lateral wall (east face) is set to an adiabatic boundary condition on both floors, following standard practice in UK residential simulation \citep{Shatat2015}, on the assumption that the adjacent dwelling maintains a similar indoor thermal environment. Second, the windows and Shading:Overhang objects on the party wall face are removed, reducing the total from 15 to 11 windows. Third, the natural ventilation objects for the party wall direction are removed (8 reduced to 6), and the Foundation:Kiva exposed perimeter fraction is reduced from 1.0 to approximately 0.75 to reflect the non-exposed party wall edge.

These modifications mean that the semi-detached archetype has less exposed surface area for solar gain, heat loss, and natural ventilation on one face --- enabling direct comparison of how built form mediates overheating risk and heating demand across the same 14-parameter design space.


\section{Parameter Space and Sampling}
\label{sec:parameters}

Fourteen building design parameters are independently sampled using continuous uniform distributions (\Cref{tab:parameters}). The ranges span from pre-1919 uninsulated stock to Passivhaus-equivalent performance, enabling the surrogate to represent the full diversity of the English detached and semi-detached housing stock.

\begin{table}[htbp]
\centering
\caption{Sampled parameters (14 continuous).}
\label{tab:parameters}
\footnotesize
\begin{tabularx}{\textwidth}{lccccX}
\toprule
\textbf{Parameter} & \textbf{Placeholder} & \textbf{Min} & \textbf{Max} & \textbf{Unit} & \textbf{Source} \\
\midrule
\multicolumn{6}{l}{\textit{Geometry (3)}} \\
Width               & \texttt{@width@}              & 6.0   & 10.0  & m       & Lawrence et al.\ (2021) \\
Length              & \texttt{@length@}             & 6.0   & 12.0  & m       & Lawrence et al.\ (2021) \\
Storey height       & \texttt{@height@}             & 2.4   & 3.0   & m       & Lawrence et al.\ (2021) \\
\midrule
\multicolumn{6}{l}{\textit{Orientation (1)}} \\
Orientation         & \texttt{@orientation@}        & 0     & 360   & deg     & Continuous; cf.\ Botti et al.\ (2022) \\
\midrule
\multicolumn{6}{l}{\textit{WWR (1)}} \\
Window-to-wall ratio & \texttt{@wwr@}               & 0.10  & 0.40  & ---     & Symonds et al.\ (2016) \\
\midrule
\multicolumn{6}{l}{\textit{Insulation (3)}} \\
Wall insulation     & \texttt{@wallInsThk@}         & 0.010 & 0.300 & m       & Part~L (2021) \\
Roof insulation     & \texttt{@roofInsThk@}         & 0.025 & 0.350 & m       & Part~L (2021) \\
Slab insulation     & \texttt{@slabInsThk@}         & 0.010 & 0.200 & m       & Part~L (2021) \\
\midrule
\multicolumn{6}{l}{\textit{Glazing (2)}} \\
Window U-factor     & \texttt{@u\_windows@}         & 0.80  & 5.00  & W/m\textsuperscript{2}K & Symonds et al.\ (2016) \\
SHGC (g-value)      & \texttt{@g\_value@}           & 0.20  & 0.70  & ---     & Lawrence et al.\ (2021) \\
\midrule
\multicolumn{6}{l}{\textit{Roof (1)}} \\
Roof absorptance    & \texttt{@roofAbs@}            & 0.30  & 0.90  & ---     & Light tiles to dark slate \\
\midrule
\multicolumn{6}{l}{\textit{Infiltration (1)}} \\
Flow coefficient    & \texttt{@flowCoeff@}          & 0.001 & 0.012 & m\textsuperscript{3}/s & Stephen (2000) \\
\midrule
\multicolumn{6}{l}{\textit{Natural ventilation (1)}} \\
Window opening factor & derived                     & 0.1   & 0.8   & ---     & Botti et al.\ (2022) \\
\midrule
\multicolumn{6}{l}{\textit{Solar shading (1)}} \\
Fixed shading depth & \texttt{@shadingDepth@}       & 0.0   & 1.0   & m       & Shading:Overhang \\
\bottomrule
\end{tabularx}
\end{table}

Eight groups of derived parameters are computed programmatically from the sampled values --- including attic zone origin, roof ridge coordinate, window coordinates from WWR, ventilation opening areas, occupant count, and daylight reference points --- yielding 34 placeholder replacements per simulation (Appendix~B). Fixed parameters including heating setpoint, cooling setpoint, internal mass, and all material properties remain constant (Appendix~C).

A Latin Hypercube Sampling design with $n = 2{,}000$ samples per archetype was generated using \texttt{scipy.stats.qmc.LatinHypercube} with a fixed random seed (seed = 42). LHS ensures uniform marginal coverage of each parameter, avoiding the clustering artefacts of purely random sampling \citep{Westermann2019}. The sample size provides approximately 143 samples per dimension, exceeding the 60--100 minimum recommended by \citet{Lomas1992}. For each sample, the batch script reads parameters from the LHS matrix, computes derived quantities, replaces 34 IDF placeholders, invokes EnergyPlus, and extracts the five target variables from the hourly CSV output. Simulations were parallelised across 8 workers, achieving approximately 5.5 seconds per run (28 minutes per archetype). All 2{,}000 simulations per archetype completed with zero failures.


\section{Target Variables}
\label{sec:targets}

Five continuous target variables are extracted from each simulation, selected to capture both regulatory overheating severity and energy performance:

\begin{enumerate}[label=(\roman*)]
    \item \textbf{$H_{e,\mathrm{worst}}$} --- Hours of Exceedance (\% of occupied hours where $\Delta T > 1$\,K), worst-case zone (TM52 Criterion~1; failure at 3\%);
    \item \textbf{$W_{e,\mathrm{max,worst}}$} --- maximum Daily Weighted Exceedance across all summer days and zones (TM52 Criterion~2; failure at $W_e > 6$);
    \item \textbf{$T_{\mathrm{op,peak}}$} --- peak operative temperature across both zones;
    \item \textbf{hours\_gt26\_night} --- nighttime hours exceeding 26\,\textdegree C in the first-floor bedroom (TM59 Criterion~B; failure at 32 hours);
    \item \textbf{heat\_kWh\_m\textsuperscript{2}} --- annual heating energy intensity (kWh/m\textsuperscript{2}$\cdot$yr).
\end{enumerate}

These are treated as continuous severity measures rather than binary pass/fail indicators, preserving modelling signal even when binary failure rates are low under the baseline climate.


\section{Surrogate Model Training}
\label{sec:training}

The dataset (2{,}000 samples $\times$ 19 columns per archetype) was partitioned into 70\% training, 15\% validation, and 15\% test subsets using a fixed random seed (random\_state = 42). The validation set was used for early stopping and learning rate scheduling during MLP training; final performance metrics were computed on the held-out test set only. For the tree-based models (RF and XGBoost), input features were left unscaled, as decision trees are invariant to monotonic transformations. For the MLP, both inputs and targets were standardised to zero mean and unit variance using separate StandardScaler instances fitted on the training set, with predictions inverse-transformed to original units for evaluation \citep{Lawrence2021,Hou2024}.

Three architectures were trained for each archetype:

\paragraph{Random Forest.} 200 trees with max\_features = $\sqrt{d}$ and default sklearn settings. Five independent models were trained, one per target variable. RF requires minimal hyperparameter tuning and provides built-in MDI feature importance \citep{Botti2022}.

\paragraph{XGBoost.} 400 estimators with learning\_rate = 0.05, max\_depth = 6, subsample = 0.8, colsample\_bytree = 0.8, min\_child\_weight = 3, L1 regularisation (reg\_alpha = 0.1) and L2 regularisation (reg\_lambda = 1.0). Five independent models were trained, one per target. Hyperparameters were manually selected based on values reported in the literature \citep{Botti2022,Shirzadi2025}; systematic hyperparameter optimisation was not performed and is noted as a limitation in \Cref{ch:discussion}.

\paragraph{MLP.} A single multi-output PyTorch neural network predicting all five targets simultaneously. The architecture consists of an input layer (14 neurons), three hidden layers ($128 \to 64 \to 32$ neurons) with Batch Normalisation and ReLU activation after each of the first two layers and Dropout (rate = 0.2) for regularisation, followed by a linear output layer (5 neurons). Training used the Adam optimiser (lr = $1 \times 10^{-3}$, weight\_decay = $1 \times 10^{-4}$) with ReduceLROnPlateau scheduling (patience = 10, factor = 0.5, min\_lr = $1 \times 10^{-5}$), MSE loss, and a batch size of 64. Early stopping with patience = 30 epochs on validation loss was applied, with a maximum of 300 epochs; the model checkpoint with the lowest validation loss was retained.

Performance was assessed on the held-out test set using $R^2$, MAE, and RMSE. Predicted-versus-actual parity plots were generated for visual assessment across the full target range.


\section{Sensitivity Analysis}
\label{sec:sensitivity}

SHAP analysis \citep{Lundberg2017} was applied to the trained XGBoost models using TreeExplainer. Global feature importance was computed as mean $|\text{SHAP value}|$ for each target. Beeswarm summary plots were generated for key targets to reveal both the magnitude and direction of each parameter's effect. RF MDI importance was computed independently to cross-validate the SHAP rankings, following \citet{BagheriEsfeh2025}. Separate SHAP analyses for the detached and semi-detached archetypes enable comparison of how built form influences the relative importance of design parameters.


% ============================================================
% CHAPTER 4: RESULTS
% ============================================================
\chapter{Results}
\label{ch:results}

% NOTE: Results below are from the detached archetype only.
% Semi-detached results to be added once simulations complete.

\section{Simulation Dataset Overview}
\label{sec:dataset-overview}

All 2{,}000 EnergyPlus simulations completed successfully. \Cref{fig:target-dist} presents the distribution of each target variable.

The Hours of Exceedance ($H_{e,\mathrm{worst}}$) distribution is heavily right-skewed and zero-inflated: 75.6\% of simulations produced $H_e = 0$, with a maximum of 0.60\% --- well below the 3\% TM52 Criterion~1 failure threshold. No parameter combination triggered a Criterion~1 failure under the London TMYx climate. This is itself a meaningful finding: sustained summertime overheating as defined by TM52 Criterion~1 is not a widespread risk for UK detached houses under the current baseline climate, even across a wide range of fabric and ventilation conditions.

In contrast, the Daily Weighted Exceedance ($W_{e,\mathrm{max,worst}}$) exhibited substantially greater variation (range 0--23, median 5.0), with 26.9\% of simulations exceeding the Criterion~2 threshold of $W_e > 6$. This indicates that while cumulative overheating is rare, short-duration peak heat stress events affect over a quarter of the parameter space --- suggesting Criterion~2 is the more relevant risk metric for this typology. Peak operative temperature ($T_{\mathrm{op,peak}}$) ranged from 27.7 to 30.6\,\textdegree C (median 28.95\,\textdegree C, $\sigma = 0.43$\,\textdegree C). Nighttime bedroom overheating (hours\_gt26\_night) was zero-inflated (64.9\% = 0) but showed a meaningful tail: 35.1\% of simulations recorded at least one nighttime hour above 26\,\textdegree C, with a maximum of 27 hours --- approaching the TM59 Criterion~B threshold. Annual heating energy intensity (heat\_kWh\_m\textsuperscript{2}) exhibited the widest dynamic range (6.8--234.6\,kWh/m\textsuperscript{2}$\cdot$yr, median 43.9).

\begin{figure}[htbp]
    \centering
    \includegraphics[width=\textwidth]{figures/target_distributions.png}
    \caption{Target variable distributions ($n = 2{,}000$).}
    \label{fig:target-dist}
\end{figure}


\section{Surrogate Model Accuracy}
\label{sec:accuracy}

\Cref{tab:results} reports test-set $R^2$ and MAE for all model--target combinations.

\begin{table}[htbp]
\centering
\caption{Test-set $R^2$ and MAE for three surrogate models across five targets.}
\label{tab:results}
\begin{tabular}{lcccccc}
\toprule
& \multicolumn{2}{c}{\textbf{Random Forest}} & \multicolumn{2}{c}{\textbf{XGBoost}} & \multicolumn{2}{c}{\textbf{MLP}} \\
\cmidrule(lr){2-3} \cmidrule(lr){4-5} \cmidrule(lr){6-7}
\textbf{Target} & $R^2$ & MAE & $R^2$ & MAE & $R^2$ & MAE \\
\midrule
$H_{e,\mathrm{worst}}$           & ---   & ---   & ---   & ---   & 0.847 & 0.013\% \\
$W_{e,\mathrm{max,worst}}$       & ---   & ---   & ---   & ---   & 0.910 & 0.82    \\
$T_{\mathrm{op,peak}}$           & 0.837 & ---   & 0.944 & ---   & 0.955 & 0.074\,\textdegree C \\
hours\_gt26\_night                & ---   & ---   & 0.743 & 0.40  & 0.725 & ---     \\
heat\_kWh\_m\textsuperscript{2}  & ---   & ---   & ---   & ---   & 0.916 & 4.76    \\
\bottomrule
\end{tabular}
\end{table}

% TODO: Fill in missing R² and MAE values from notebook output

The MLP consistently achieved the highest $R^2$ across four of five targets: $T_{\mathrm{op,peak}}$ ($R^2 = 0.955$, MAE = 0.074\,\textdegree C), heat\_kWh\_m\textsuperscript{2} ($R^2 = 0.916$, MAE = 4.76\,kWh/m\textsuperscript{2}), $W_{e,\mathrm{max,worst}}$ ($R^2 = 0.910$, MAE = 0.82), and $H_{e,\mathrm{worst}}$ ($R^2 = 0.847$, MAE = 0.013\%). XGBoost achieved the highest $R^2$ for hours\_gt26\_night ($R^2 = 0.743$, MAE = 0.40 hours), marginally outperforming MLP ($R^2 = 0.725$). RF ranked third across all targets, with $R^2$ values 5--15 points below MLP.

$T_{\mathrm{op,peak}}$ was best-predicted across all models (MLP $R^2 = 0.955$; XGBoost 0.944; RF 0.837), reflecting the smooth, near-linear relationship between fabric parameters and summer peak temperature. Nighttime bedroom overheating was the most difficult target (best $R^2 = 0.743$), attributable to its zero-inflated distribution and the discontinuous response near the 26\,\textdegree C threshold.

\begin{figure}[htbp]
    \centering
    \includegraphics[width=\textwidth]{figures/r2_comparison.png}
    \caption{$R^2$ per target variable for MLP, RF, and XGBoost.}
    \label{fig:r2-comparison}
\end{figure}


\section{Parity Analysis}
\label{sec:parity}

\Cref{fig:parity} presents parity plots comparing predictions against EnergyPlus values on the test set. For $T_{\mathrm{op,peak}}$ and heat\_kWh\_m\textsuperscript{2}, both XGBoost and MLP produce tight clusters along the 1:1 diagonal. For $W_{e,\mathrm{max,worst}}$, scatter increases above $W_e > 10$, with MLP achieving better fit ($R^2 = 0.910$ vs 0.886). For $H_{e,\mathrm{worst}}$, points cluster near the origin, with greater scatter for non-zero values --- reflecting the challenge of predicting a zero-inflated target. The hours\_gt26\_night plots show the most scatter between 1--10 hours, though both models correctly distinguish zero from non-zero cases.

\begin{figure}[htbp]
    \centering
    \includegraphics[width=\textwidth]{figures/parity_plots.png}
    \caption{Parity plots: XGBoost vs MLP for all five targets.}
    \label{fig:parity}
\end{figure}


\section{Neural Network Training Behaviour}
\label{sec:training-behaviour}

\Cref{fig:loss-curves} shows MLP training and validation loss curves over 160 epochs. Validation loss converges by approximately epoch 40 and remains stable, with no evidence of overfitting. Early stopping was employed with patience = 20 epochs.

\begin{figure}[htbp]
    \centering
    \includegraphics[width=\textwidth]{figures/loss_curves.png}
    \caption{MLP training and validation loss (MSE) over 160 epochs.}
    \label{fig:loss-curves}
\end{figure}


\section{Feature Importance and Sensitivity Analysis}
\label{sec:feature-importance}

\subsection{SHAP Global Feature Importance}

\Cref{fig:shap-global} presents mean $|\text{SHAP value}|$ for the XGBoost model across all five targets.

For overheating metrics ($H_{e,\mathrm{worst}}$, $W_{e,\mathrm{max,worst}}$, $T_{\mathrm{op,peak}}$), WWR and g-value are consistently the two most important features. Fixed shading depth ranks third; roof insulation fourth. Orientation has minimal SHAP importance under the London TMYx climate, likely because the typical meteorological year averages out directional solar gain effects.

For nighttime bedroom overheating (hours\_gt26\_night), window opening factor is overwhelmingly dominant, followed by g-value, wall insulation, and WWR --- corroborating \citet{Botti2022}. Nighttime overheating is primarily a ventilation problem; wall insulation's secondary importance reflects thermal mass effects.

For heating energy (heat\_kWh\_m\textsuperscript{2}), WWR ranks first, followed by wof, wall insulation, g-value, and length. WWR's dominance reflects the dual role of glazing: solar gain reduction and conductive heat loss increase.

\begin{figure}[htbp]
    \centering
    \includegraphics[width=\textwidth]{figures/shap_global.png}
    \caption{SHAP global feature importance per target.}
    \label{fig:shap-global}
\end{figure}

\subsection{SHAP Summary Plot ($H_{e,\mathrm{worst}}$)}

\Cref{fig:shap-beeswarm} shows the SHAP beeswarm plot for $H_{e,\mathrm{worst}}$. High WWR values increase overheating severity in a non-linear, threshold-like pattern. High g-value pushes $H_{e,\mathrm{worst}}$ upward; deep overhangs reduce it. Roof insulation shows a non-monotonic pattern: both thin and very thick insulation increase $H_{e,\mathrm{worst}}$ --- thin insulation permits direct roof heat conduction, while thick insulation traps internal gains, consistent with the overheating--insulation paradox documented by \citet{Petrou2021}.

\begin{figure}[htbp]
    \centering
    \includegraphics[width=0.85\textwidth]{figures/shap_beeswarm.png}
    \caption{SHAP beeswarm summary plot for $H_{e,\mathrm{worst}}$.}
    \label{fig:shap-beeswarm}
\end{figure}

\subsection{Random Forest Feature Importance}

\Cref{fig:rf-mdi} shows RF MDI feature importance, which broadly corroborates the SHAP rankings. The agreement between SHAP (model-agnostic) and MDI (RF-specific) provides cross-validation of the identified parameter sensitivities.

\begin{figure}[htbp]
    \centering
    \includegraphics[width=\textwidth]{figures/rf_mdi.png}
    \caption{RF MDI feature importance per target.}
    \label{fig:rf-mdi}
\end{figure}


\section{Summary of Key Findings}
\label{sec:results-summary}

First, MLP achieved the highest overall accuracy ($R^2 > 0.9$ for three of five targets). Second, under the London TMYx climate, no parameter combination triggered TM52 Criterion~1 failure, but 26.9\% exceeded Criterion~2 --- indicating peak-day heat stress is the more relevant metric. Third, SHAP identified WWR and g-value as primary daytime overheating drivers and wof as the primary nighttime driver. Fourth, the overheating--insulation paradox is empirically confirmed: roof insulation shows non-monotonic effects on $H_{e,\mathrm{worst}}$.


% ============================================================
% CHAPTER 5: DISCUSSION (TODO)
% ============================================================
\chapter{Discussion}
\label{ch:discussion}

% TODO: Write Discussion chapter
% Sections to include:
% \section{Model Performance in Context}
\section{Limitations}
\label{sec:limitations}
% TODO: Write limitations section
% \section{Implications for Practice}


% ============================================================
% CHAPTER 6: CONCLUSION (TODO)
% ============================================================
\chapter{Conclusion}
\label{ch:conclusion}

% TODO: Write Conclusion chapter


% ============================================================
% REFERENCES
% ============================================================
\bibliography{references}


% ============================================================
% APPENDICES (TODO)
% ============================================================
\appendix
% \chapter{Construction Details}
% \chapter{Derived Parameter Calculations}
% \chapter{Fixed Parameters}

\end{document}
